\chapter{Physical Activity}
\index{Physical Activity}

\medskip
\noindent\textit{This chapter is for general education. Increase activity gradually and ask a healthcare professional about safe limits if you have a medical condition.}

\section*{What it is}
Physical activity is any movement that uses energy, including walking, cycling, housework, play, sport, work, and exercises that strengthen muscles or improve balance. Regular activity supports heart, bone, muscle, metabolic, and mental health.

\section*{How much is helpful}
Adults should aim for regular moderate or vigorous activity across the week, with muscle-strengthening activity on at least 2 days when able. Children, older adults, pregnant people, and people with disability may need tailored advice. Any activity is better than none; sitting for long periods can be interrupted with short movement breaks.

\section*{Starting safely}
Choose activities you enjoy and begin at a manageable level. Warm up if helpful, use suitable footwear or equipment, drink according to conditions and exertion, and increase duration or intensity slowly. Strength, balance, and flexibility work can complement aerobic activity. Stop and seek advice for persistent pain or symptoms that are unusual for you.

\section*{When to Seek Medical Care}
Stop exercising and seek urgent help for chest pressure, severe breathlessness, fainting, new neurological symptoms, or sudden severe pain. Ask for individual advice before strenuous activity if you have heart or lung disease, are recovering from illness or surgery, are pregnant with complications, or have another condition affecting exercise tolerance.

\section*{Sources and further reading}
\begin{itemize}
\item World Health Organization. \textit{Physical activity}. \url{https://www.who.int/news-room/fact-sheets/detail/physical-activity}
\item Centers for Disease Control and Prevention. \textit{Benefits of Physical Activity}. \url{https://www.cdc.gov/physical-activity-basics/benefits/index.html}
\end{itemize}
